\chapter{The removal of agency by server owners and reclaiming player power}

\begin{center}
    Author: matcha\_tea. Editor: retrooper. Single-blind peer reviewed. % clearly indicate non-peer reviewed articles! use \footnote{...} for more information
\end{center}

\begin{flushright}
    IMSC: j(2)-ma-8

    DOI: \doilink{10.5281/zenodo.23702589}
\end{flushright}

% ============================================================================

\section*{Abstract}
Throughout the history of Minecraft servers, there is one system that has remained the same: leadership. At the top is the server’s owner, the person who has the final say in any decision made; next are the server’s admins, a group of people whom the owner trusts to act on their behalf; finally, there are the moderators, who enforce the server’s rules and maintain conduct in the day-to-day running of the server. I call this the Traditional structure. This paper argues that this structure is fundamentally flawed by exploring the core function and purpose of a Minecraft server, which is to create \emph{satisfaction} for players on the server. Two key terms are outlined: \emph{agency} and \emph{power}, which are important for understanding why the Traditional structure has and continues to fail players by examining the issue from the perspective of the Player. This paper considers the individual \emph{player} and the lone decision-maker – the owner – as classes in a hierarchy to better understand the inter-class interactions, removing individual anomalies in behaviour and allowing readers to understand the problem at a higher level. Finally, this paper proposes a structure where Players (as a class) have no “need” for the Owner, as decisions would only be made by Players.

\section{Introduction}
The rigid Owner-Player hierarchy found in Traditional structures has created the most (monetarily) successful servers in the world (e.g., Hypixel, Mineplex, and Donut SMP). Both Raid Farms and Minecraft Bloodlines (MCB), the most studied and influential servers in the field of Minecraft Server Theory (MST), have this rigid top-down structure of power. It can appear logical, even ideal, to have a few people at the top dictating the direction of the server and viewing players as numbers and statistics needing improvement.

Such hierarchical systems do not exist in isolation, but in the context of everything around their creators. Just like the most popular Minecraft servers, the most well-known real-life corporations also follow this kind of hierarchy – Amazon, Nestle, and Tesla to name a few.

I will argue in this paper that this hierarchical structure fails to achieve its goal, even when one thinks solely about higher and higher player counts. To understand why, one must also reframe one’s thinking of the meaning of “success” for a server, which I argue should be in terms of \emph{player satisfaction} instead of other numerical metrics. Following this, one must also understand that the Traditional structure fails to meet this goal and is also the reason why the vast majority of servers fail after the Two-Week Phase. I will discuss an alternative way of thinking about player satisfaction, and by extension, the success of a server. Moreover, having understood why the Traditional structure fails and how a different view should be adopted, I will discuss how the same structure deals with player suggestions and change. I will argue that the unilateral nature of the Owner means that Players are not satisfied, resulting in the eventual collapse in feelings of excitement for the server’s potential. Finally, I will lay out a better structure of leadership which holds Player satisfaction at its core and would yield a more “successful” server.

In this paper, all mentions of “server” can be assumed to mean a \emph{Survival Multiplayer (SMP)} server. An SMP is best defined in its components, which includes but is not limited to:\begin{itemize}
    \item Survival gamemode
    \item Multiplayer
    \item Focused on player-to-player interaction
    \item An open and expansive world for players to do as they wish
\end{itemize}

\subsection{Use of ``satisfaction'' in this paper}
The word \emph{satisfaction} is used throughout this paper instead of \emph{fun} as, ironically, what gives players \emph{satisfaction} may not always be \emph{fun}. Take the example of fighting the Ender Dragon. It is a difficult fight for the typical player, and one may die several times during the fight. This is by no means an activity that can objectively and easily be considered \emph{fun}, but it brings about satisfaction in that the player continues to engage in the activity and has fun in the end. One can consider satisfaction to be fun in hindsight, where even if a player struggled during a difficult part of the game, as long as they came out of it for the better, they can think \emph{in hindsight} that they \emph{did} have fun.

\section{Existing structures and hierarchies}
The current mainstream leadership structure is roughly similar to the following:
\begin{quote}
    (In descending order of power)

    Owner > Admin > Moderator > Player
\end{quote}
As explained in the next section, this structure is inefficient in achieving player satisfaction. However, we must first discuss \emph{how} this structure works in practice and \emph{why} it sometimes creates a successful server.

The Owner\footnote{\emph{Owner} and other terms are capitalised as I am talking about the institution and the class of person, not specifically the person who embodies the role. This distinction is important as the things I talk about are \emph{inherent} to the role and will not change depending on \emph{who} is in that specific office.} holds the most power in a server and all decisions must be made with their approval. As they hold control of the Minecraft server (and usually also the Discord server that they might have), only \emph{they} can make changes. Even if it is a decision that a majority of the leadership team (that is, the Admins and possibly the Moderators, together with the Owner) supports, only the Owner is able to actually install a mod or change the game difficulty, for example. This means they have complete autocratic control over all things, even if they are amicable and adhere to the majority rule. Thus, I think it is less pertinent to discuss the roles of the Admins and Moderators, as their power is ultimately given by the Owner.

If a real-life example can aid your understanding, take the United Kingdom – the King is the \emph{de facto} head of state. Any law passed by both the Houses of Commons and Lords must be signed off by the King so that it may become the official law of the land. Without his rubber stamp, any law coming onto his desk will not be in effect, even if the entire country supports it.\footnote{Parliament can, using specific laws, force legislation to be approved.}

This is completely and undoubtedly dangerous. The high concentration of power, the selection of the next leader by the current leader, and the leader's failure to understand the experiences of ordinary people compound into a dangerous slope toward detachment from reality.

Why must an Owner, who is not elected, not chosen, and cannot be held accountable through any means (as they have the ultimate say and can completely disregard all dissent), control the entire server? Or really, in concrete terms – how does this ensure player satisfaction?

In a typical SMP server founded between friends, one of them will be the Owner who has likely initiated the idea, owning the physical server that the world is hosted on. Initially, the Owner may decide to ask their Players what kinds of mods and mechanics they would want present on the server. Thereafter, some players may request a change, and others may disagree with the requested change. The Owner acts as a mediator in these cases to make a quick decision based on their own opinion to move the conversation forwards, as they simply want to have fun with their friends. Here, maintaining a strict democratic decision-making process is not necessary as the goal is not to boost player satisfaction and increase the regular player count; instead, the owner is part of a group of friends who use their server as a \emph{medium} to interact and spend time together.

In other words, disagreements in terms of gameplay mechanics can be taken less seriously due to the server’s goal, so having efficient and unilateral decision-making is good. Unilateral decision-making may also be more beneficial if all players already know each other before starting to play – which means discussion and compromise is more likely to be fruitful.

One may raise the issue that the server’s goal might have been influenced or introduced by the players at the beginning of the server, which means that even if the owner unilaterally makes executive decisions, it still follows the intention set by the players at the beginning.

This is true; however, we must remember that goals and ideas shift over time. A player who suggests, for example, the Copper Rails mod to be installed on the server as they believe it will increase the uses of copper may later find, after using it for a while, that it actually decreases the incentive to mine for gold and iron. Here, their opinion on copper rails may change, but yet, they have no power to change the fact that the mod is installed.

In order for a server to have broad appeal, ensuring player satisfaction \emph{and} increasing its regular player count, two things must be considered: player \emph{agency} and player \emph{power}.

I shall define \emph{agency} as follows:
\begin{quote}
    Agency can be thought of as a player’s inward feeling about whether they have the power to bring about change effectively on a server—which is a standard definition.
\end{quote}
Power, however, has nuance to its definition in this context, which I shall define as follows:
\begin{quote}
    Typically, power is thought of as something which is direct, physical (or mental), or radical. However, power can come in many different ways. This can be power to speak up to a person with more power (e.g., the Owner), power to influence other players, power to enact practical change, and power to simply have one’s opinion heard.
\end{quote}
Let us first explore a simplified scenario involving a new player who initially enjoys a server but ultimately decides to leave.

First, the player is intrigued by something on the server. If we take the example of Raid Farms, this may be its famous wheat currency. They join the server and they are engaged with this mechanic. As time passes, they notice more things about the new mechanic that they previously hadn’t, continuing to make them more and more enthused.

Then, the player inevitably finds a mechanic that they are not the most pleased with. Sticking with our example, the enchanting system has been a large topic of interest and discussion. At this point, the player feels a lot of agency to be \emph{able} to do this as they have a positive view of the server overall, but not yet a lot of power in seeing real change or feedback.

Next, our newly joined but already dedicated player wants to change their new favourite server for the better. They make a suggestion in the server and find that many agree with them. Here, their feeling of agency continues to rise as their power to bring about change is affirmed. They also begin to feel more power as they seem to have invited many people who agree with them, rallying for the change.

The Owner notices this, agrees, and listens to their feedback. We see the peak of both agency and power here: agency is felt by those who feel strongly about this issue when the highest authority on the server affirms that their concerns are valid, and power is felt because they managed to influence this person, who is seen as a role model.

Finally, when the change is implemented, the player might feel unsatisfied, especially if the implementation does not meet their expectations. The change may not have been implemented perfectly, or it may even just require one or two more tweaks. However, at this point, the player’s agency is completely shattered as they feel that they are not heard, and thus, did not bring about the change they desired.

Not all players decide to leave immediately after this, but it is very likely that some might. We should also remember that our dedicated new joiner would not leave at this point, as they feel a sense of \emph{loyalty} to the server, but the vast majority of players view a server as a product to be consumed, not an institution or organisation deserving protection, advancement, and dedication.\footnote{I recognise that this is an unsupported claim.}

In terms of looking at agency and power for player satisfaction, the breaking point occurred when the Owner, being this perceived bastion of power and being the ultimate decision-maker, went against the wishes of the players. The players’ individual belief that the Owner knows what is best for the server, and what steps to take in order to reach their goal, has been shattered as they \emph{themselves} put the Owner on this pedestal.

One may argue that the fault is not in having an Owner but in the beliefs of the Players. However, this ignores the context in which the Players exist – one where the Owner places themselves in this position of power and status, almost asking to be admired for their apparent wisdom. Players simply follow this expectation that the Owner must be looked up to – creating the situation we find ourselves in.

\section{Pleasing the players}
Now that we’ve discussed what typical mainstream servers have as their hierarchy of power, and the feelings of agency and power individual players have, we shall shift focus to the Players as a group.

Say that you are the owner of a server – how do you balance your own vision for the server from before it started with the assumed visions of the players? How can you make sure the decisions demanded by the players are recognised without allowing them to make choices that (in your view) harm the server? Is it possible to increase the popularity of your server without assimilating with the ever-growing number of servers who all have one, stubborn, goal in their mind – increasing player count?

In order to answer these questions, you must understand the two weights in play in any multiplayer server – the view of the players and the view of the server owner.\footnote{In some servers, this may also be the view of the admins, or even moderators. In any case, I refer to the person, or group of people, who have power to make decisions about the server.} If the two views are visualised as a two-circle Venn diagram, the area of their intersection can be thought of as the amount of flexibility an owner has in making decisions – how much of their own circle they can pick ideas out of. This intersection also represents the views that both the Owner and the Players have in common.

However, these are \emph{weights}, meaning that they also possess another property – their physical mass – or, how important they are perceived to be. This perception can be held by the server owner regarding players, meaning how much of the players’ opinions are taken into account when deciding server mechanics, for example. This can also be a perception held by the players themselves – whether they feel they have a voice in making decisions.

Four possible situations arise from these two factors which I shall deal with independently:\begin{enumerate}
    \item Large intersection of views\footnote{That is, of the server owner and the players.} + High weight of player opinion
    \item Large intersection of views + Low weight of player opinion
    \item Low intersection of views + High weight of player opinion
    \item Low intersection of views + Low weight of player opinion
\end{enumerate}
The first case is the optimal case in a hierarchical server – the server owner is satisfied that their direction for the server lines up with their player base’s opinion, and they are open to taking the players’ views in consideration when making decisions.

One may argue that the second term is a direct cause of the first term. In other words, because the owner has a large number of views in common with their players, it is much easier for the owner to accept them as peers. Although this view seems straightforward and natural, the two terms are completely independent from each other. Take the example of a server owner whose direction for their server is popular among the player base but who does not actively listen to them, instead making decisions independently – it is simply coincidental that the owner’s actions align with the players’ wants.

The above situation would be representative of the second case – a lot in common, but the owner acts independently. It is also important to note that these terms are fluid – an owner who has a lot in common with their players at the establishment of the server may not hold the same values or vision for the server after a few months of operation; similarly, an owner who is receptive to feedback at the start may later find their player base acting rashly, and thus, decrease the weight on player opinion when making decisions.

The third case is of a humble and open owner – although their vision is unaligned with their players’, they are still open to feedback and advice. These owners listen to their players attentively and are usually the ones who are thought of most highly by the players and considered “one of the people”, creating an equal environment where all views are respected.

One may argue here that the third case is the optimum. Firstly, despite the differences in opinion, the owner continues to take the players’ perspectives into account; thus, this must be the case where the most respect is given – even more than the first case, as the respect in that case may be fragile. Further, the owner being amicable with their player base, and the players being amicable with the owner, must mean that a more positive community is built.

To explain why both arguments fail – or why both positions remain rather \emph{inefficient} for the goal of player satisfaction – we must consider them in relation to the goal of a server.

What is the purpose of a Minecraft server? Or even, what \emph{is} a Minecraft server?

I will define a server (that is, not \emph{only} an SMP, e.g., anarchy, geopol) for the purposes of the following sections as follows:
\begin{quote}
    A multiplayer environment where people involved in the environment, be it as a moderator or player, work towards a goal that they individually have. This may be a shared goal, a formally recognised and rigid goal, or even an informal goal such as “fun”. An “environment” could refer to a Survival world with the Nether and End locked, or a Creative server with building plots.
\end{quote}
Ultimately, servers have a loose goal – and for an SMP, which this paper deals with – this is \emph{fun}. The exact definition of fun is not important here, but the fact that this subjective goal remains desired is of importance.\footnote{Note that \emph{fun} and \emph{satisfaction} are not the same idea. \emph{Fun} is a broad and generally undefined word used to describe a “vibe” players get when they are feeling \emph{satisfied} on a server. Because it is unspecific and relating to an \emph{emotion}, I separate these.}

There are a few things that must be understood to realise the flawed nature of one figure ruling over an entire server.

Firstly, recognise that \emph{players} play on a server. Even if the owner plays on the server, they can also be thought of as a player. The subjective goal of fun needs to be agreed upon by the majority of the players in order to maximise the amount of fun that any particular individual may have. This is usually done through discussion between players either within the server environment itself (in-game) or through external forums, such as Discord.

Secondly, understand that the players \emph{themselves} will know best about what \emph{they} want. Even when given a private, honest space to share their opinions, some players will still feel anxious about dissenting against the Owner, whom they perceive as “all-knowing”.

Further, no amount of oversight or observation by the owner may be able to represent \emph{all} player opinions. Even if the owner plays on the server regularly and engages with the community as a normal player, they will still have a skewed or biased view of the server. There is also the question of differential treatment by players towards the owner when they are playing, which will also exaggerate the positive aspects of the players and diminish the negative ones.

One can argue that if the owner plays on an “alt account”, they will be able to get a true experience on the server; however, this will also produce a biased result depending on the owner’s own play style, influencing whom they interact with.

Thirdly, many players often put forward suggestions for changing a server by discussing them among themselves, then bringing those suggestions forward to the owner – usually through more influential or well-known players who may have sway over the owner. However, this is an inefficient way of creating change.

Fourthly, even if the owner always does approve of all suggestions, there may be disagreement in the implementation of the suggestions. An owner will always have a skewed view of their server, as mentioned above, due to their status and role in the power hierarchy. This may have the effect of creating more harm than good, as the players could feel that their desires are not being fulfilled – but at the same time, the owner can develop a perception that their work is not being valued and that implementing player suggestions is a waste of time.\footnote{Of course, not every owner will think the same. However, for the purposes of this paper, I shall use a generalisation of the worst-case-scenario where applicable.} Consequently, this creates a negative feedback loop that causes the area of intersection between the Owner and the Players to shrink, leading to the opinion of the Players being given less weight, until the point where the two classes are completely polarised and players begin to leave.

Having considered these points, one can still find a commonality between them – which is that all decisions could be made by the owner. One may suggest that this results in quick, efficient decision-making and less bureaucracy, allowing players to experience server changes more quickly and accelerating the feedback loop.\footnote{The feedback loop is the following: problem → solution identified → change → player feedback → (repeat)} This may also mean that players have a higher likelihood (due to how frequently opportunities to show agency arise) of feeling satisfaction. All processes have trade-offs, and this example is one that we must deal with.

In order to please players and restore their power, inefficiency is a necessary evil. As we’ve established, only players know best what they want in a server, and any other person or group of people making decisions \emph{for} them in the name of “knowing better” is \emph{more} inefficient as it will not deliver a satisfactory result and will require more iterations of the feedback loop – this is compared to the often tedious and long-drawn nature of discussion and collective decision making, but it almost always guarantees a more satisfactory result.

The fact that players collectively make these decisions also means that discussions can be more civil as every person has equal amounts of power, sharing a unifying goal that is constructive – compared to the unifying, but unhelpfully critical, goal that players may have against the owner.

One could construct a strong argument that setting up a new server and establishing new leadership can provide an entirely clean slate for a more stable form of leadership. One fault in this argument is that the players’ emotional attachment to the world remains overlooked. This means a brand-new world must be created. Players must re-establish relationships, create new cultural landmarks, and invent new in-jokes about the server. This process is physically exhausting, mentally exhausting, and most importantly: it decreases satisfaction monumentally.

Further, I believe the revolution (so to say) of the Players against the Owner can deliver even more satisfaction than resetting everything. Server cultures stay in place, relationships remain stable (if not reinforced), and the world stays the same. If, eventually, the Players do overthrow the Owner, a new vibrant culture of daily discussion – which players know will at least not be a waste of time – can drive players’ motivation to play and participate on the server. This is in contrast to there being an issue on a server with a Traditional structure, where a player may try to create change before ultimately giving up.

\section{Reclaiming player power}
The players’ power can, and must be, \emph{reclaimed}. The system of operating a server where this is achieved is the \emph{Co-operative} system.

Due to the already established monopoly of the Owner over the Players, it is essential that the Players unite against the Owner. There must be no doubt between players that what they are against is an (almost) immovable force who has control over \emph{their} server and \emph{their} community. The Players must realise that as long as they have an owner, their feelings of agency will be capped at a certain level due to the presence of an immovable obstacle to gaining ultimate collective power.

This is where \emph{celebrity players} come in – these people have the most agency (as they already know that they are popular amongst other players), and the most power (as they are also recognised \emph{personally} by the owner), which allows them to influence both the classes effectively. The influence they perform leads to increased feelings of agency, which in turn also boosts their own satisfaction – which can lead to a positive feedback loop, allowing them to influence the owner even more.

I envision this to be a long, tedious, and very difficult process: when you try to dispel the myth of “efficiency” and “wisdom” of the Owner to a recipient who has only ever known the status-quo, you will meet great resistance. However, once much of the Players have been educated about these matters, they should constantly and without hesitation repeat the same, commonly agreed-upon message of disapproval to show their unity against the Owner. The Players shall force the Owner to be removed and immediately set up concrete infrastructure within the Discord server (or other means of communication) to establish the following as a minimum:\footnote{Please keep in mind that this should not be taken as a fully planned out manifesto – this paper simply lays the foundation for its necessity.}\begin{enumerate}
    \item An open, free, and respectful space for debate and discussion
    \item Direct voting on server mechanics, which can be initiated by any player
    \item An Executive Council chosen by lottery to implement approved changes
\end{enumerate}
The first is essential for any Player-owned server, as the basis of the entire concept is for players to be able to discuss amongst themselves and enact change directly to other members of the same class; a safe space for all players to participate in this process is absolutely necessary. This space must be inclusive of all and do its best to give equal weight and time to all voices, no matter how famous, influential, or veteran. However, at the same time, members of the Executive Council must not be able to participate in debate, and they must be politically neutral because of their role, which requires them to be disinterested and unbiased parties whose job is simply to \emph{carry out} the agreed-upon wishes of the Players.

This will also cement the fact that agency and satisfaction are the primary goals of a co-operative server. If players feel that they are able to drive change (and if they truly can drive change), they will feel like they have \emph{agency}. Similarly, if there is a lower barrier to enjoying oneself on the server, satisfaction will come about more easily.

The second of the three-legged stool allows for the results of these debates to be carried out. Direct voting means that every player’s vote is equal and meaningful. Decisions by referendum also avoid the problem that ultimately rises from using a representative system, which gives certain players an elevated position, mimicking the Owner. The latter part of this point enshrines the equal right and power of every player to be able to create change. In terms of agency and power, this reaffirms both.

This right also creates opportunities for abuse and trolling, as a newly joined troll has the exact same power as a two-year veteran. However, this right is essential to maintain a “no exceptions” culture when it comes to having the right to say something is wrong, and to change it, starting from oneself. To address this risk, a simple majority vote could be conducted, serving as a preliminary check before continuing to debate a proposed issue or solution.

The final point ensures that the debates held were not pointless and that a team of regular players – not a dictatorial entity like the Owner – will carry out the practical steps needed implement a new change. This Executive Council is chosen by lottery, meaning that it is not disproportionately filled with self-sorting players who may have a certain bias or obsession with the server’s operations, making them \emph{interested} parties and likely ineffective at being unbiased when carrying out the necessary, practical work. This also serves as a safeguard against corruption, as the players who bring a proposal to the Council cannot influence it without knowing \emph{who} is on the Council.

A random chance of being on the Council also means that a culture of equality is established. Because serving on it is normalised, players begin to feel like part of a political culture where everyone \emph{is} truly \emph{seen} and \emph{treated} as equal.

\section{Caveat for this system}
It is important to note that, just like many server proposals, one that focuses on restoring agency and power may not always be the best for the players. Take the example of a Vanilla UHC server – there are a set of rules, some mods, and everything else is set in stone. Not only are the Players’ opinions not heard or cared about, but they are not given any weight either. The established culture within the server is that “it is what it is”, and making changes is irrelevant. The prevailing view is that a server built on survival, competition, PvP, and defending your base does not require players to get involved with \emph{how} it is managed.

Although total disregard for player feedback remains rare, servers with concrete and specific goals face a significant challenge. It is difficult to justify introducing player-led decision-making on these servers, even regarding peripheral issues where the players might actually make better decisions.

\section{Conclusion}
In this paper, we’ve discussed existing hierarchical structures of leadership, player satisfaction regarding agency and power, and the intersections of Player and Owner views. We also assessed how much weight these views hold for the Owner and provided a template for how Players can take back ownership of their servers.

\section{Acknowledgements}
Thanks to Harlow for giving me feedback on the first draft of this paper. I would also like to thank the two anonymous reviewers and retrooper, the editor, who all gave extensive feedback on the manuscript.

\section{Declarations}
No conflict of interest and no AI use have been declared.